\bookchapter{第五十六章：打磨與發佈——測試、除錯與上架 Google Play}

\chapterintro{App 做出來了，然後呢？這一章回答“從能跑”到“能交付”的三件事：怎麼高效地除錯找 bug、怎麼寫測試防迴歸、怎麼把 App 送上 Google Play 讓全世界下載。上架不只是儀式感——它是把一個作品變成產品的分水嶺。}

\booksection{一、除錯三板斧：Logcat、斷點、UI 檢查器}

\begin{center}
\begin{tabularx}{\linewidth}{|X|X|X|}
\hline
\textbf{工具} & \textbf{幹什麼} & \textbf{怎麼用} \tabularnewline
\hline
Logcat & 看系統/自己打的日誌 & 程式碼裡 Log.d("TAG","值=\$x")，面板裡搜 TAG 過濾 \tabularnewline
斷點除錯 & 讓程式停在某行，看變數值 & 行號旁點一下出紅點，Debug 模式執行，逐步走 \tabularnewline
Layout Inspector & 看執行中的介面結構 & 選單 Tools Layout Inspector，實時看元件樹 \tabularnewline
\hline
\end{tabularx}
\end{center}

\begin{center}\small\color{gray}除錯三板斧\end{center}

\begin{tcbinfo}{小貼士：報錯怎麼讀？}
崩潰後 Logcat 裡會有紅色堆疊（Stacktrace），從下往上找 第一個自己包名開頭的行（比如 com.example.yourapp），那行基本就是案發現場。先看“Exception 型別”（空指標？型別轉換？），再看“at ... 你的檔案:行號”。這招能解決 80\% 的崩潰。

\end{tcbinfo}

\booksection{二、測試：給程式碼上保險}

手測靠不住——改了這個功能，另一個悄悄壞了。測試就是“後悔藥”兼“防彈衣”。安卓測試分三層，新手先吃透第一層：

\begin{center}
\begin{tabularx}{\linewidth}{|X|X|X|X|}
\hline
\textbf{層級} & \textbf{測什麼} & \textbf{放哪裡} & \textbf{速度} \tabularnewline
\hline
單元測試 & 純邏輯（計算、校驗） & app/src/test/ & 秒級 \tabularnewline
儀器測試 & 真實裝置上的介面行為 & app/src/androidTest/ & 分鐘級 \tabularnewline
手動測試 & 真人感受、體驗 & —— & 永遠需要 \tabularnewline
\hline
\end{tabularx}
\end{center}

\begin{center}\small\color{gray}測試三層\end{center}

\begin{lstlisting}[language=Kotlin]
// 把"金額字串轉數字"這種邏輯抽成純函式
fun parseAmount(input: String): Double? {
    val v = input.trim().toDoubleOrNull()
    return if (v != null && v > 0) v else null
}

// 測試它（JUnit4）
@Test
fun emptyInputReturnsNull() {
    assertEquals(null, parseAmount(""))
}
@Test
fun negativeReturnsNull() {
    assertEquals(null, parseAmount("-5"))
}
@Test
fun normalValueParses() {
    assertEquals(12.5, parseAmount("12.5"))
}
\end{lstlisting}

記帳本里的金額校驗、統計計算，都該有單元測試罩著。跑測試：右鍵 test 目錄 Run。綠色全過，放心改程式碼。

\booksection{三、簽名打包：給 App 蓋上“印章”}

安卓要求所有安裝包必須簽名——就像蓋公章，證明“這個包是我出的，沒被篡改”。除錯時 Android Studio 自動用 debug 簽名；發佈要用自己的 release 簽名。流程：

\begin{enumerate}
  \item 選單 Build Generate Signed Bundle / APK；
  \item 建立 keystore 檔案（金鑰庫）——設定密碼並妥善保管；
  \item 選擇 release + AAB 格式，構建完成。
\end{enumerate}

\textbf{避坑提醒：keystore 丟了 = App 丟了}\par

上架後更新 App 必須用同一個簽名，否則商店會拒絕更新。keystore 檔案和密碼一旦丟失，你只能換個新應用重新上架——老使用者全都接不上。務必把 keystore 和密碼備份到安全的地方（雲盤加密 + 記事本），這是安卓開發者最昂貴的教訓之一。

\booksection{四、上架 Google Play 全流程}

註冊開發者帳號

準備商店素材

上傳 AAB

填資料安全表

稽核（數小時\textasciitilde{}數天）

上架

\begin{center}\small\color{gray}圖 6：Google Play 上架六步走\end{center}

\begin{itemize}
  \item 註冊：一次性 25 美元註冊費，需能用的 Google 帳號；
  \item 素材：應用圖示（512×512）、截圖（至少 2 張）、功能描述、隱私政策連結（收集資料必填）；
  \item 資料安全表：如實申報收集了哪些資料——不聯網的記帳本可以"零收集"；
  \item 稽核：新帳號首次上架常被人工稽核，數小時到數天不等，被打回就按意見改。
\end{itemize}

國內使用者注意：Google Play 需要特殊網路環境，且中國大陸地區分發通常走國內商店（華為、小米、應用寶等）。上架國內商店的流程類似，但各商店有自己的稽核口味——同一個 AAB/APK 可以多店分發。作為練手，先走通一家就贏了。

\booksection{五、商店素材怎麼做：第一印象就是下載量}

商店頁是使用者“逛店”時看到的櫥窗。三個素材決定第一印象：

\begin{itemize}
  \item 圖示：512×512，別用截圖湊合。沒設計能力就用 AI 工具生成一個簡潔幾何圖形，在手機桌面上縮小後還能一眼認出才算合格；
  \item 截圖：至少 2 張，最好 4-6 張，展示最核心的功能場景；每張配一句人話文案（"一秒記一筆"比"高效的記帳解決方案"好十倍）；
  \item 描述：第一句說清"這是什麼 + 給誰用"，別堆砌"最專業、最強"——商店搜尋權重主要看關鍵詞和下載轉化，真誠比誇張有效。
\end{itemize}

\begin{tcbinfo}{小貼士：A/B 看看同行}
上架前逛一圈商店裡同類 App 的前十名，截圖對比：他們圖示什麼風格？截圖講什麼賣點？標題怎麼寫？模仿不是抄襲——學套路、換自己的內容，這是所有產品冷啟動最快的一課。

\end{tcbinfo}

\booksection{五、效能最佳化與崩潰上報小清單}

App 上架前，這兩件事值得各花半小時：

\booksubsection{效能自查四條}

\begin{enumerate}
  \item 列表用 LazyColumn：長列表千萬別用 Column 硬排——懶載入是安卓效能的第一課；
  \item 圖片別硬塞原圖：用 Coil 圖片庫自動縮放快取，一張 5MB 原圖足以卡爆低端機；
  \item 避免在主執行緒乾重活：資料庫、網路、大檔案都在協程裡做；
  \item 用 Profiler 量一量：Android Studio 自帶的效能分析器能看記憶體、CPU、耗電——先測量，再最佳化，別憑感覺。
\end{enumerate}

\booksubsection{崩潰上報}

你測試時沒崩，不代表使用者手機不崩。上架前接入 Firebase Crashlytics（谷歌官方，免費）：幾行配置，使用者崩潰時自動上報堆疊，你在後臺就能看到“哪個機型、哪個版本、崩在哪一行”。沒有崩潰監控的上架，等於閉著眼睛開車。

\begin{tcbpractice}{來動手：把記帳本"交付"一次}
(1) 給 parseAmount 等邏輯補 3 個單元測試；(2) 用 Logcat 親手抓一個自己製造的崩潰並修好；(3) 生成一個 release 簽名的 APK，發給朋友安裝（這才是“作品”的儀式感）；(4) 有餘力就註冊開發者帳號，把記帳本真正提交上架。前三條今晚就能做完。

\end{tcbpractice}

\begin{tcbsummary}{本章小結}
\begin{itemize}
  \item 除錯三板斧：Logcat 看日誌、斷點看變數、Layout Inspector 看介面
  \item 讀崩潰：從堆疊下往上找自己包名的第一行
  \item 測試三層：單元（秒級，先學這個）/儀器/手動
  \item 發佈必須簽名：keystore 和密碼備份好，丟了 App 就"換人"了
  \item 上架六步：註冊 素材 上傳 AAB 資料安全表 稽核 上架
\end{itemize}

\end{tcbsummary}
